\subsection{Метод узловых потенциалов}
%%%%%%%%%%%%%%%%%%%%%%%%%%%%%%%%%%%%%%%%%%%%%%%%%%%%
%%%%%%%%%%%%%%%%%%%%%%%%%%%%%%%%%%%%%%%%%%%%%%%%%%%%1
 \z[{\begin{wrapfigure}[4]{r}{0.25\textwidth}%
		\vspace{-20pt}%
		\begin{flushright}%
 			\includegraphics[width=0.25\textwidth]{3.7.1}%
		\end{flushright}%
	\end{wrapfigure}%
	}]%
{В схеме на рисунке все элементы можно считать идеальными. Значения $\emf$ источников и сопротивлений резисторов на схеме указаны. Определите величину и направление тока через амперметр.
\vspace{2em}
}
%
{$I = 1\,\text{А}$, вниз}

%%%%%%%%%%%%%%%%%%%%%%%%%%%%%%%%%%%%%%%%%%%%%%%%%%%%2
 \z[{\begin{wrapfigure}[5]{r}{0.25\textwidth}%
		\vspace{-30pt}%
		\begin{flushright}%
 			\includegraphics[width=0.25\textwidth]{3.7.2}%
		\end{flushright}%
	\end{wrapfigure}%
	}]%
{Электрическая цепь, схема которой изображена на рисунке, состоит из двух батареек с $\emf_1 = 5 \ \text{В}$ и $\emf_2 = 2 \ \text{В}$ 
и пяти одинаковых вольтметров. Найдите показания каждого из вольтметров. Сопротивления батареек много меньше сопротивлений вольтметров.
\vspace{2em}
}
%
{$U_1 = U_2 = \frac{2}{7}(\emf_1 + \emf_2) = 2\,\text{В}$; $U_3 = \frac{1}{7}(3\emf_1 - 4\emf_2) = 1\,\text{В}$; \par $U_4 = \frac{1}{7}(4\emf_1 - 3\emf_2) = 2\,\text{В}$; $U_5 = \frac{3}{7}(\emf_1 + \emf_2) = 3\,\text{В}$}

%%%%%%%%%%%%%%%%%%%%%%%%%%%%%%%%%%%%%%%%%%%%%%%%%%%%3
 \z[{\begin{wrapfigure}[4]{r}{0.25\textwidth}%
		\vspace{-40pt}%
		\begin{flushright}%
 			\includegraphics[width=0.25\textwidth]{3.7.3}%
		\end{flushright}%
	\end{wrapfigure}%
	}]%
{Найдите силу тока через перемычку $ab$ в схеме, изображённой на рисунке. 
Сопротивлениями перемычки, подводящих проводов и внутренним сопротивлением батареи пренебречь.
\vspace{2em}
}
%
{$I = \frac{\emf}{7R}$}


%%%%%%%%%%%%%%%%%%%%%%%%%%%%%%%%%%%%%%%%%%%%%%%%%%%%4
 \z[{\begin{wrapfigure}[2]{r}{0.25\textwidth}%
		\vspace{-35pt}%
		\begin{flushright}%
 			\includegraphics[width=0.25\textwidth]{3.7.4}%
		\end{flushright}%
	\end{wrapfigure}%
	}]%
{При каком сопротивлении $R_x$ ток через резистор $R$ равен нулю? Параметры цепи указаны на схеме.
\vspace{2em}
}
%
{$I_1 = \frac{\emf}{R}$; $I_2 = \frac{3\emf}{R}$; $I_3 = \frac{2\emf}{R}$}

%%%%%%%%%%%%%%%%%%%%%%%%%%%%%%%%%%%%%%%%%%%%%%%%%%%%
 \z[{\begin{wrapfigure}[3]{r}{0.25\textwidth}%
		\vspace{-20pt}%
		\begin{flushright}%
 			\includegraphics[width=0.25\textwidth]{3.7.5}%
		\end{flushright}%
	\end{wrapfigure}%
	}]%
{Определите разность потенциалов между точками $a$ и $b$ в схеме, изображённой на рисунке. Параметры схемы указаны.
\vspace{2em}
}
%
{$\frac{18\emf}{11}$}

%%%%%%%%%%%%%%%%%%%%%%%%%%%%%%%%%%%%%%%%%%%%%%%%%%%%6
 \z[{\begin{wrapfigure}[4]{r}{0.25\textwidth}%
		\vspace{-10pt}%
		\begin{flushright}%
 			\includegraphics[width=0.25\textwidth]{3.7.6}%
		\end{flushright}%
	\end{wrapfigure}%
	}]%
{В электрической цепи, схема которой показана на рисунке, все элементы идеальные, их параметры указаны.
\lpu Найти ток через ключ $K$ с указанием направления после замыкания ключа. \lpuend
\lpu Найти отношение напряжений между точками $A$ и $B$ после и до замыкания ключа. \lpuend
\vspace{2em}
}
%
{\puan $\frac{\emf}{8R}$, налево;  \quad \puan $\frac{U_2}{U_1} = \frac{9}{10}$}

%%%%%%%%%%%%%%%%%%%%%%%%%%%%%%%%%%%%%%%%%%%%%%%%%%%%
 \z[{\begin{wrapfigure}[4]{r}{0.45\textwidth}%
		\vspace{-35pt}%
		\begin{flushright}%
 			\includegraphics[width=0.45\textwidth]{3.7.7}%
		\end{flushright}%
	\end{wrapfigure}%
	}]%
{В схеме, изображённой на рисунке, определить показания амперметров и вольтметра. Приборы считать идеальными.}
%
{$I_1 = 0$; \quad $I_2 = \frac{\emf}{8R}$; \quad $V = 0$}



